\chapter{Methodology / project development}
\label{chap:methodology}

This compact methodology follows the paper structure while keeping the thesis
evidence that matters: the experimental contract, the development lessons that
justify the final route, the detailed frozen-prior formulas, and the final
prior-anchored residual GMM head. The longer derivation files
\texttt{methodology.tex}, \texttt{prior\_detail\_try78.tex},
\texttt{prior\_detail\_try79.tex} and \texttt{prior\_detail\_try80.tex}
remain in the repository as source-level mathematical notes; the compiled text
below keeps the same formulas in a shorter paper-style form.

\section{Dataset and evaluation contract}

Each sample is a complete \(513\times513\) CKM with
\SI{1}{m}\(\times\)\SI{1}{m} pixels. The UAV transmitter is always
horizontally centred in the map, but its antenna height \(h_{\mathrm{tx}}\)
varies by sample and is one of the main physical conditioning variables. Every
non-building pixel is an active receiver at
\(h_{\mathrm{rx}}=\SI{1.5}{m}\). Building pixels are not receiver locations and
are removed from the loss and metrics.

The antenna-height distribution is skewed rather than uniform. Reading the
stored \texttt{uav\_height} field from the HDF5 database gives an empirical
range of \SIrange{11.61}{477.54}{m}. The closest tested nominal beta-mixture
sampler is approximately
\(h_{\mathrm{design}}\approx 10+490X\), with
\(X\sim0.75\,\mathrm{Beta}(3,23)+0.25\,\mathrm{Beta}(4,4)\), but the realized
database distribution is represented by the \SI{3}{m} empirical density below.
The value \SI{477.54}{m} is the largest sampled height, not a hard endpoint of
the approximate beta law; drawing no samples above \SI{478}{m} is statistically
normal because the expected count above that height is only 0.52:
\begin{equation}
  \hat f_{\Delta h}(h) = d_k,\quad
  d_k=\frac{n_k}{N\Delta h},\quad
  h\in[10+3k,\,13+3k),\quad
  \Delta h=\SI{3}{m},\quad N=16180 ,
  \label{eq:height_empirical_histogram}
\end{equation}
where \(n_k\) is the number of stored heights in bin \(k\). Display curves
linearly interpolate between bin centres. The all-database counts are 3972
samples below \SI{50}{m}, 8456 from \SIrange{50}{150}{m}, 2490 from
\SIrange{150}{300}{m}, and 1262 above \SI{300}{m}; in the training cities the
same bins contain 2676, 5654, 1644, and only 866 samples. Thus the
high-altitude bins later showing lower RMSE are easier despite having fewer
training samples, mainly because they contain fewer blocked links and weaker
spread tails.

\begin{table}[htbp]
\centering
\scriptsize
\caption{City-holdout assignment used by the distribution-first and final neural models.}
\label{tab:split_city_assignment}
\begin{tabularx}{\textwidth}{lrrX}
\toprule
Split & Samples & Cities & City names \\
\midrule
Train & 10840 & 37 & Abidjan, Abu Dhabi, Alexandria, Antigua Guatemala,
Athens, Bangalore, Bangkok, Bath, Berlin, Bratislava, Brugge, Buenos Aires,
Canberra, Chefchaouen, Dalian, Hanoi, Ho Chi Minh City, Hyderabad, Kinshasa,
Leuven, Ljubljana, Luanda, Lyon, Manaus, Minsk, Nagpur, Paris,
Rio de Janeiro, Santiago, Shanghai, Siem Reap, St.~Louis, Thane, Toledo,
Valparaiso, Venice, Vienna. \\
Validation & 2750 & 15 & Bilbao, Chiang Mai, Cleveland, Dhaka, Dubrovnik,
Fortaleza, Marrakesh, Mexico City, Munich, Nazareth, Nice, Salvador, Segovia,
Tunis, Victoria. \\
Test & 2590 & 14 & Barcelona, Beijing, Carcassonne, Copenhagen, Halifax,
Jaipur, Johannesburg, Key West, Kuala Lumpur, Osaka, Phuket, Pune, Surat,
Vancouver. \\
\bottomrule
\end{tabularx}
\end{table}

The receiver-valid mask is
\begin{equation}
  m_{\mathrm{ground}}(i,j)=
  \begin{cases}
    1,& \mathrm{topology}(i,j)=0,\\
    0,& \mathrm{otherwise},
  \end{cases}
  \label{eq:ground_mask}
\end{equation}
so all reported pixel errors are pooled over the non-building receivers:
\begin{equation}
  \mathrm{RMSE}
  =
  \sqrt{
  \frac{\sum_{s,i,j} m_s(i,j)(\hat y_s(i,j)-y_s(i,j))^2}
       {\sum_{s,i,j} m_s(i,j)}
  }.
  \label{eq:rmse}
\end{equation}
The split is city-level, not random-pixel or random-crop. This is essential:
the model is evaluated on unseen urban layouts rather than on other samples
from cities already seen during training.

\section{Development lessons retained from the attempt history}

The methodology changed in exposition, not in the final experimental contract.
The long chronological version recorded many intermediate attempts; the
compact version keeps only the evolution that explains why the final method is
structured as it is.

\begin{enumerate}
\setlength\itemsep{0.15em}
\item Direct U-Net and cGAN baselines validated the aligned image-to-image
pipeline, but adversarial visual realism did not solve physical error.
\item The ground-mask correction made building pixels invalid for training and
reporting; older numbers became historical trend markers rather than a single
leaderboard.
\item PMNet/PMHHNet-style residual learning showed that height conditioning,
LoS/NLoS diagnosis and physical channels help, but point regression still
collapsed hard NLoS tails.
\item Distribution-first GMM models proved that the NLoS error was not
inevitable: representing the target distribution fixed much of the collapse
even without explicit physical priors.
\item The final path-loss and spread priors showed that large parts of the
problem are better handled by train-only calibrated estimators.
\item The final model combines both lessons: frozen priors explain stable
structure and a residual GMM head learns bounded local corrections.
\end{enumerate}

Thus the prior is useful but not magical. It is not external knowledge that
solves the task by itself; it is a strong, frozen, auditable estimator of the
part of the map that is easier to calibrate than to relearn. LoS path loss,
NLoS path loss, delay spread and angular spread are treated differently because
their empirical distributions are different.

\section{Final pipeline}
\label{sec:priors_overview}

\begin{figure}[htbp]
\centering
\resizebox{0.98\textwidth}{!}{%
\begin{tikzpicture}[
  font=\scriptsize,
  box/.style={draw, thick, rounded corners=2pt, rectangle, minimum width=2.75cm, minimum height=0.75cm, align=center},
  input/.style={box, fill=gray!15},
  prior/.style={box, fill=blue!8},
  cal/.style={box, fill=green!10},
  model/.style={box, fill=orange!12},
  arr/.style={-{Stealth[length=2.1mm]}, thick, rounded corners=3pt}
]
\node[input] (support) at (0,0) {Topology + masks\\centred UAV\\variable \(h_{\mathrm{tx}}\)};
\node[prior] (los) at (3.4,1.35) {LoS PL\\FSPL + two-ray\\height-bin curves};
\node[prior] (nlos) at (3.4,0) {NLoS PL\\COST/A2G envelope\\\(\mathbf{x}_{78}\) OLS};
\node[prior] (spread) at (3.4,-1.35) {DS/AS\\\(\log(1+y)\) raw prior\\\(\mathbf{x}_{79}\) ridge};
\node[cal] (stack) at (6.9,0) {Frozen prior maps\\PL, DS, AS};
\node[model] (net) at (10.2,0) {Shared U-Net\\FiLM height conditioning\\residual GMM heads};
\node[input, fill=orange!18] (out) at (13.4,0) {Continuous PL, DS, AS};
\draw[arr] (support.east) -- ++(0.45,0) |- (los.west);
\draw[arr] (support.east) -- (nlos.west);
\draw[arr] (support.east) -- ++(0.45,0) |- (spread.west);
\draw[arr] (los.east) -- ([yshift=0.28cm]stack.west);
\draw[arr] (nlos.east) -- (stack.west);
\draw[arr] (spread.east) -- ([yshift=-0.28cm]stack.west);
\draw[arr] (stack.east) -- (net.west);
\draw[arr] (net.east) -- (out.west);
\end{tikzpicture}}
\caption{Final prior-anchored CKM pipeline. Frozen height-aware priors are
computed first and then used as input channels for the residual GMM model.}
\label{fig:pipeline}
\end{figure}

The model receives nine aligned channels: topology, LoS mask, NLoS mask,
ground mask, combined path-loss prior, LoS path-loss prior, NLoS path-loss
prior, delay-spread prior and angular-spread prior. The scalar
\(h_{\mathrm{tx}}\) is encoded with sinusoidal features and injected through
FiLM modulation \cite{perez2018film}. The network therefore learns a residual
around already height-aware priors instead of inferring all height effects from
pixels alone.

\section{Path-loss prior}
\label{sec:try78_detail}

The path-loss prior (Try 78 in the implementation records) has two branches.
The LoS branch is a coherent two-ray model plus calibrated height/radius
corrections. For each receiver pixel, with \(d_{2D}\) measured from the centred
transmitter,
\begin{equation}
\begin{aligned}
d_{\mathrm{LoS}}&=\sqrt{d_{2D}^2+(h_{\mathrm{tx}}-h_{\mathrm{rx}})^2},\\
d_{\mathrm{ref}}&=\sqrt{d_{2D}^2+(h_{\mathrm{tx}}+h_{\mathrm{rx}})^2}.
\end{aligned}
\label{eq:d_los_ref_compact}
\end{equation}
The free-space term is
\begin{equation}
\mathrm{FSPL}
=32.45+20\log_{10}\!\left(\frac{d_{\mathrm{LoS}}}{1000}\right)
+20\log_{10}(f_{\mathrm{MHz}}).
\label{eq:fspl}
\end{equation}
The coherent reflection correction is
\begin{equation}
C_{\mathrm{2ray}}
=-20\log_{10}\!\left|
1+\rho(h_{\mathrm{tx}})\frac{d_{\mathrm{LoS}}}{d_{\mathrm{ref}}}
e^{-j(k(d_{\mathrm{ref}}-d_{\mathrm{LoS}})+\phi(h_{\mathrm{tx}}))}
\right|.
\label{eq:c2ray}
\end{equation}
The deployed LoS prior is
\begin{equation}
\mathrm{PL}_{\mathrm{LoS}}
=\mathrm{FSPL}+C_{\mathrm{2ray}}+b(h_{\mathrm{tx}})
+r(d_{2D},h_{\mathrm{tx}}).
\label{eq:pl_los}
\end{equation}
Here \(\rho\) and \(\phi\) are effective reflection amplitude and phase,
fitted in \SI{5}{m} height bins, smoothed across bins and interpolated to the
actual UAV height. The term \(b\) removes a remaining height-dependent offset
and \(r\) is a smoothed radial residual table. These extra terms are outside
the textbook two-ray formula; they make the prior match the CKM dataset while
keeping the ring-like LoS structure explicit and auditable.

The NLoS branch begins with a conservative raw blocked-link prior. A
COST231-shaped urban term is
\begin{equation}
\begin{aligned}
\mathrm{PL}_{C}
={}&46.3+33.9\log_{10}(f_{\mathrm{MHz}})
-13.82\log_{10}(h_{\mathrm{tx}})-a(h_{\mathrm{rx}})\\
&+\left(44.9-6.55\log_{10}(h_{\mathrm{tx}})\right)
\log_{10}(d_{\mathrm{km}})+3,
\end{aligned}
\label{eq:cost231}
\end{equation}
and an elevation-aware A2G-NLoS envelope is
\begin{equation}
\mathrm{PL}_{\mathrm{A2G,NLoS}}
=\lambda_0(h_{\mathrm{tx}},f)+\beta_{\mathrm{NLoS}}
+A_{\mathrm{NLoS}}\exp\!\left(-\frac{90-\theta^\circ}{\tau_{\mathrm{NLoS}}}\right),
\label{eq:a2g_nlos}
\end{equation}
with \(\lambda_0=20\log_{10}(4\pi h_{\mathrm{tx}}f/c)\). The raw NLoS prior is
\begin{equation}
P_0=\max(\mathrm{PL}_{C},\mathrm{PL}_{\mathrm{A2G,NLoS}}).
\label{eq:nlos_raw}
\end{equation}
The final NLoS estimate is not this raw formula. It is a regime-wise linear
calibration over morphology and height features:
\begin{equation}
\mathrm{PL}_{\mathrm{NLoS}}
=\operatorname{clip}(\mathbf{x}_{78}^{\top}\mathbf{w}_{78,r},20,180),
\label{eq:pl_nlos_final}
\end{equation}
where
\begin{equation}
\begin{aligned}
\mathbf{x}_{78}=[&
P_0^2,P_0,\log(1+d_{2D}),\rho_{15},\rho_{41},h_{15},h_{41},\\
&\rho_{41}h_{41},\rho_{41}\log(1+d_{2D}),n_{15},n_{41},
s,\theta,n_{41}\theta,1]^{\top}.
\end{aligned}
\label{eq:x78}
\end{equation}
\(\rho_k\) is local building density in a \(k\times k\) window, \(h_k\) is a
local building-height statistic, \(n_k\) is local NLoS support, \(s\) is a
shadow/blocker proxy and \(\theta\) is normalized elevation. The interaction
terms let dense tall areas differ from sparse tall areas and let obstruction
matter differently at different ranges and elevations. The coefficient vector
\(\mathbf{w}_{78,r}\) is fitted on training cities only, with \(r\) selected
from topology, LoS/NLoS and antenna-height regime fallbacks.

The deployed path-loss prior is selected by the geometric visibility mask:
\begin{equation}
\mathrm{PL}_{\mathrm{prior}}
=m_{\mathrm{LoS}}\mathrm{PL}_{\mathrm{LoS}}
+(1-m_{\mathrm{LoS}})\mathrm{PL}_{\mathrm{NLoS}}.
\label{eq:pl_hybrid}
\end{equation}

\section{Delay and angular spread priors}
\label{sec:try79_detail}

The spread priors (Try 79 in the implementation records) are fitted in
\(\log(1+y)\) space:
\begin{equation}
z=\log(1+y),
\qquad
\hat y=\exp(\hat z)-1.
\label{eq:log_transform_79}
\end{equation}
This is necessary because delay spread and angular spread have a LoS spike near
zero plus NLoS-heavy tails. The raw log prior is
\begin{equation}
\begin{aligned}
z^{(0)}
={}&\log(1+b^{\mathrm{nat}}_q)
+b_{\mathrm{top}}
+a_1\log(1+d_{2D})+a_2\theta_{\mathrm{inv}}\\
&+a_3\rho_{41}+a_4h_{41}+a_5n_{41}
+a_6n_{41}\theta_{\mathrm{inv}},
\end{aligned}
\label{eq:raw_spread}
\end{equation}
where \(q\) is the LoS/NLoS state, \(b^{\mathrm{nat}}_q\) sets a native LoS or
NLoS scale, \(b_{\mathrm{top}}\) is a topology-class shift,
\(\theta_{\mathrm{inv}}\) is the low-elevation complement, and the morphology
terms describe local density, height and NLoS support.

The calibrated spread feature vector is
\begin{equation}
\begin{aligned}
\mathbf{x}_{79}=[&
(z^{(0)})^2,z^{(0)},\log(1+d_{2D}),\theta_{\mathrm{norm}},
\theta_{\mathrm{inv}},h,h^2,\\
&\rho_{15},\rho_{41},h_{15},h_{41},n_{15},n_{41},
n_{41}\log(1+d_{2D}),\\
&\rho_{41}\theta_{\mathrm{inv}},b_{41},1,c_{41},t_{41},
\theta_{\mathrm{norm}}\rho_{41},\\
&c_{41}\theta_{\mathrm{inv}},t_{41}\rho_{41},t_{41}n_{41}]^{\top}.
\end{aligned}
\label{eq:x79}
\end{equation}
Here \(h\) is normalized UAV height, \(b_{41}\) is blocker depth,
\(c_{41}\) is transmitter clearance over nearby buildings, and \(t_{41}\) is
the fraction of locally taller buildings. These terms make antenna height a
first-class feature: the prior can distinguish a UAV above the local skyline
from one near rooftop level.

For each metric and regime, ridge calibration solves
\begin{equation}
\hat{\mathbf{w}}_{79,r}
=
\left(\mathbf{X}^{\top}\mathbf{X}
+\lambda\mathbf{D}^{\top}\mathbf{D}\right)^{-1}
\mathbf{X}^{\top}\mathbf{z}.
\label{eq:ridge}
\end{equation}
The spread dictionary contains 114 fitted keys, but this is not a single
Cartesian product. For each spread metric, \(6\) topology classes times \(2\)
visibility states times \(3\) height bins give \(36\) exact keys. The
implementation adds \(12\) height-collapsed topology/visibility fallbacks,
\(6\) topology-wide fallbacks, \(2\) global visibility fallbacks and \(1\)
fully global fallback, for \(57\) keys per metric and \(114\) across DS and AS.

\section{Prior-anchored residual GMM-head model}
\label{sec:try80_detail}

The final model (Try 80 in file names and JSON logs) predicts residual
distributions around the frozen priors. For target
\(t\in\{\mathrm{PL},\mathrm{DS},\mathrm{AS}\}\), visibility regime \(q\) and
pixel \(x\), the head predicts mixture weights and bounded residual offsets:
\begin{equation}
\mu_{t,k}(x)
=
p_t(x)+c_{t,q(x)}\tanh r_{t,k}(x),
\label{eq:try80_mu}
\end{equation}
where \(p_t(x)\) is the frozen prior, \(c_{t,q}\) is the residual cap for the
target/regime and \(r_{t,k}\) is the learned raw residual. The deterministic
map exported for metrics is the mixture expectation
\begin{equation}
\hat y_t(x)
=
\sum_{k=1}^{K}\pi_{t,k}(x)\,\mu_{t,k}(x),
\qquad K=3 \text{ in the final residual head.}
\label{eq:try80_point}
\end{equation}
The GMM head therefore has two roles: it represents residual uncertainty during
training and it produces a continuous pixel-wise prediction at inference.

\begin{figure}[htbp]
\centering
\resizebox{0.98\textwidth}{!}{%
\begin{tikzpicture}[
  font=\scriptsize,
  box/.style={draw, thick, rounded corners=2pt, rectangle, minimum width=2.9cm, minimum height=0.72cm, align=center},
  inp/.style={box, fill=gray!14},
  enc/.style={box, fill=blue!8},
  head/.style={box, fill=orange!12},
  arr/.style={-{Stealth[length=2.0mm]}, thick, rounded corners=3pt}
]
\node[inp] (x) at (0,0) {9 map channels\\+ \(h_{\mathrm{tx}}\) code};
\node[enc] (stem) at (3.2,0) {Shared U-Net\\GroupNorm + FiLM};
\node[head] (pl) at (6.6,1.1) {PL residual GMM\\\(\pi,\mu,\sigma\)};
\node[head] (ds) at (6.6,0) {DS residual GMM\\\(\pi,\mu,\sigma\)};
\node[head] (as) at (6.6,-1.1) {AS residual GMM\\\(\pi,\mu,\sigma\)};
\node[inp, fill=green!10] (out) at (10.1,0) {Mixture expectation\\dense PL/DS/AS maps};
\draw[arr] (x.east) -- (stem.west);
\draw[arr] (stem.east) -- ++(0.3,0) |- (pl.west);
\draw[arr] (stem.east) -- (ds.west);
\draw[arr] (stem.east) -- ++(0.3,0) |- (as.west);
\draw[arr] (pl.east) -- ++(0.25,0) |- (out.west);
\draw[arr] (ds.east) -- (out.west);
\draw[arr] (as.east) -- ++(0.25,0) |- (out.west);
\end{tikzpicture}}
\caption{Final residual GMM-head structure. The shared network predicts
task-specific residual distributions; the pixel-wise output is the expectation
of each anchored mixture.}
\label{fig:try80_arch}
\end{figure}

\section{Objective and training settings}
\label{sec:try80_objective}

The full objective can include mixture NLL, empirical-distribution KL,
moment matching, prior anchor, prior guard, outlier budget and deterministic
RMSE/MAE terms:
\begin{equation}
\begin{aligned}
\mathcal{L}
=\sum_t w_t(&
\lambda_{\mathrm{NLL}}\mathcal{L}_{\mathrm{NLL},t}
+\lambda_{\mathrm{KL}}\mathcal{L}_{\mathrm{KL},t}
+\lambda_{\mathrm{mom}}\mathcal{L}_{\mathrm{mom},t}\\
&+\lambda_{\mathrm{anchor}}\mathcal{L}_{\mathrm{anchor},t}
+\lambda_{\mathrm{guard}}\mathcal{L}_{\mathrm{guard},t}
+\lambda_{\mathrm{out}}\mathcal{L}_{\mathrm{out},t}\\
&+\lambda_{\mathrm{RMSE}}\mathcal{L}_{\mathrm{RMSE},t}
+\lambda_{\mathrm{MAE}}\mathcal{L}_{\mathrm{MAE},t}).
\end{aligned}
\label{eq:try80_total_loss}
\end{equation}
All terms are masked by \(m_{\mathrm{ground}}\). The final selected checkpoint
is the result of a deterministic RMSE-dominant fine-tuning stage, so the
probabilistic head is part of the architecture but the final claim is an RMSE
improvement over frozen priors, not a calibrated uncertainty claim.

\begin{table}[htbp]
\centering
\caption{Active loss and training settings in the final model family.}
\label{tab:try80_loss_settings}
\scriptsize
\begin{tabularx}{\textwidth}{lXX}
\toprule
Setting & Residual/GMM stage & Final fine-tuning \\
\midrule
NLL / KL / moment / outlier & 1.00 / 0.25 / 0.10 / 0.02 & 0 / 0 / 0 / 0 \\
RMSE / MAE & 1.50 / 0.10 & 1.00 / 0.05 \\
Anchor / guard & 0.05 / 0.10 & 0.01 / 0.02 \\
Task weights & PL 1.0, DS 1.0, AS 1.0 & PL 1.0, DS 0.25, AS 0.25 \\
Learning rate & \(2\times10^{-4}\) & \(5\times10^{-5}\) \\
Backbone & 128 base channels, 160-dimensional condition vector & same architecture \\
\bottomrule
\end{tabularx}
\end{table}

The practical sequence was: wide residual model, relaxed residual caps,
large-capacity residual model, and final RMSE-dominant fine-tuning. Older
PMHHNet training devices such as SWA, EMA, CutMix, label noise and many
auxiliary losses were intentionally not part of the final selected stage. The
final design is simpler: train-only priors first, then a bounded residual model
whose gains are measured against those priors on unseen cities.
